\documentclass[12pt]{article}
\usepackage[utf8]{inputenc}
\usepackage[T1]{fontenc}
\usepackage{hyperref}
\title{Datooga Language Resources for K-12 Curriculum Development}
\author{Compiled by Claude Sonnet 5 (claude-sonnet-5) for Andr\'as Kornai}
\date{}
\begin{document}
\maketitle
\section{Introduction}
Datooga (ISO 639-3 tcc) is a Nilotic language of Tanzania. Search results describe the Barabaig dialect (around Mount Hanang) as the best known; Wikipedia notes that Barabaiga and Gisamjanga have been written and that work on Asimjeeg is under way. Dialect differences are large. Speaker estimates in the sources differ widely (Joshua Project about 16,000--17,000 for the language entries seen; American Bible Society cites about 200,000 people). Very little curricular material turned up in search.
\section{Primary Educational Resources}
No primers, textbooks or school programmes for Datooga were found in search results. The American Bible Society reports that Datooga had no written form until 1985, that literacy is below one percent, and that most Datooga do not speak Swahili.
\section{Linguistic and Reference Materials}
No dictionary or grammar was found in the search results. The Wikipedia language page gives dialect and orthography background: \url{https://EN.wikipedia.com/wiki/ISO_639:tcc}
\section{Religious and Community Resources}
The American Bible Society describes a new Datooga Bible translation: \url{https://news.americanbible.org/article/new-translation-for-tanzania}

Joshua Project profiles: \url{https://joshuaproject.net/people_groups/18375/TZ} and \url{https://joshuaproject.net/people_groups/15286/TZ}
\section{Digital and Audio Resources}
Global Recordings Network lists audio and audio-visual programmes in Barabaig (Datooga): \url{https://globalrecordings.net/ja/language/tcc}
\section{Curriculum Development Context}
The language is not known to be taught in schools, and the base of written material is thin. A curriculum would first need orthography consolidation across dialects, new primers and a lexicon. Swahili is the likely bridging language.
\end{document}